\documentclass[36pt,extrafontsizes,article,oneside,final]{memoir}
%http://mirror.las.iastate.edu/tex-archive/macros/latex/contrib/memoir/memman.pdf
% see page 46 for font size options table 3.9
\usepackage[utf8]{inputenc}
\usepackage[ngerman]{babel}
\usepackage{amsmath}
\usepackage{amsfonts}
\usepackage{amssymb}
\usepackage{graphicx}
\usepackage{calc}
\usepackage{ccicons}
%\usepackage{lipsum}
\usepackage{float}
\usepackage{picture}
\usepackage[absolute]{textpos}
\usepackage{rotating}
\usepackage{geometry}
\geometry{a0paper,left=1cm,right=1cm,top=5cm,bottom=11cm}
%\geometry{showframe=true} 
\usepackage{layout}

\headheight 3.5cm              %% put this outside
\headsep 6.5cm                 %% put this outside


\usepackage{background}
\backgroundsetup{
	scale=1,
	color=black,
	opacity=1,
	angle=0,
	contents={%
		\includegraphics[width=\paperwidth,height=\paperheight]{parts/background-image2}
	}%
}


\usepackage{chronosys}

\usepackage{xcolor}

\definecolor{faculty}{RGB}{0,61,143}

%change to iftrue to turn debug boxes on
\iffalse
	\newcommand{\debugBox}[1]{\fbox{#1}}
\else
	\newcommand{\debugBox}[1]{#1}
\fi

%change font to Modern Computer
\renewcommand{\familydefault}{\sfdefault}

\makeoddhead{plain}{}{\color{white} \textbf{Ergebnisse eines Seminars im Bachelorstudiengang
5-Fachsemester \\ }}{\color{white} Bennet Bleßmann \\ Henning Sachmann}
\makeoddfoot{plain}{\ccLogo \ccAttribution \ccShareAlike}
{\tiny Dieses Material steht unter der Creative-Commons-Lizenz Namensnennung - Weitergabe unter gleichen Bedingungen 4.0 International.}
{\tiny Um eine Kopie dieser Lizenz zu sehen, besuchen Sie http://creativecommons.org/licenses/by-sa/4.0/.}

\author{Bennet Bleßmann, Henning Sachmann}
\title{Was kann schief gehen bei der Entwicklung von algorithmischen Entscheidungsfindungsprozessen?}
\begin{document}

\begin{textblock*}{0.9\textwidth}(0.1\textwidth+2.25cm,22.0cm)
%chronology
\begin{turn}{270}
\begin{minipage}[t][0.85\textwidth]{\textheight-325pt}
\startchronology[
	startyear=0,
	stopyear=15,
	dates=false, % don't display beginning and end date 
	height=100pt,
	color=faculty]
	
\setupchronoevent{date=false,
	markdepth=40pt,
	textstyle=\Large
}
\setupchronoperiode{bottomdepth=10pt,
	dates=false,				
	topheight=90pt,
	textdepth=-120pt,
	textstyle=\Large
}

\chronoperiode[textdepth=350pt]{0}{3}{\rotatebox{90}{\colorbox{blue}{\color{white}Anforderung}}}
\chronoevent{1}{\rotatebox{90}{Kundenanforderung}}
\chronoevent{2}{\rotatebox{90}{Machbarkeitsprüfung}}

\chronoperiode{3}{6}{\rotatebox{90}{\colorbox{red}{Entwurf}}}
\chronoevent[markdepth=-120pt]{4}{\rotatebox{90}{Bestimmung der benötigten Daten}}
\chronoevent[markdepth=-120pt]{5}{\rotatebox{90}{Wahl/Entwurf des Algorithmus}}

\chronoperiode[textdepth=950pt]{6}{8}{\rotatebox{90}{\colorbox{cyan}{Entwicklung/Umsetzung des Entwurfs}}}

\chronoperiode[textdepth=-150pt]{8}{11}{\rotatebox{90}{\colorbox{purple}{\color{white}Training}}}
\chronoevent[markdepth=-120pt]{9}{\rotatebox{90}{Sammeln/Filtern der Trainingsdaten}}
\chronoevent[markdepth=-120pt]{10}{\rotatebox{90}{Training mit den Trainingsdaten}}

\chronoperiode[textdepth=550pt]{11}{15}{\rotatebox{90}{\colorbox{yellow}{Wartung/Anwendung}}}

\stopchronology

\end{minipage}
\end{turn}
\end{textblock*}

\begin{minipage}[c][\textheight-50pt][t]{0.50\textwidth}
\begin{center}
\makebox(900,460){
\parbox{900pt}{
\begin{center}
\begin{Spacing}{2}
\textbf{
%left of logo
\huge
Mögliche Probleme bei der Anwendung von Deep-Learning
}
\end{Spacing}
\end{center}
}}

\makebox(480,320){
%left of Anforderung
}

\colorbox{gray!20}{
\makebox(900,480){
%left of Entwurf
\large
\scalebox{1.2}{
\parbox{850pt}{
\begin{itemize}[$\bullet$]
\item Falsches Modell
\item Fehlende Sicherheitsüberlegung
\end{itemize}
}}}}

\makebox(110,270){
%left of Entwicklung
}

\colorbox{gray!20}{
\makebox(900,460){
%left of Training
\large
\scalebox{1.2}{
\parbox{850pt}{
\begin{itemize}[$\bullet$]
\item Unzureichende Abbildung der \\ Realität durch die Trainingsdaten
\item Verunreinigte Daten
\item Irrelevante Daten
\end{itemize}
}}}}

\makebox(500,430){
%left of Wartung
}

\colorbox{gray!20}{
\makebox(900,520){
\parbox{850pt}{
%bottom left

%\large

Lösungen für diese Probleme muss auch die Gesellschaft finden,
häufig überwiegen aktuell die finanziellen Interessen.
Durch transparenteren und nachvollziehbareren Code hätte
man z.B. den Abgasskandal verhindern können. 
Wir müssen uns bewusst machen, dass diese Entscheidungen nicht von Wenigen hinter verschlossenen Türen getroffen werden sollten,
da alle Mitglieder der Gesellschaft davon,
z.B. durch den Schufa-Score, betroffen sind.

}}}

\end{center}
\end{minipage}
\begin{minipage}[c][\textheight-50pt][t]{0.45\textwidth}

\makebox(900,320){
%ontop of logo
}

\colorbox{gray!20}{
\makebox(900,480){
%right of Anforderung
\large
\scalebox{1.2}{
\parbox{850pt}{
\begin{itemize}[$\bullet$]
\item Unklarer Kundenwunsch
\item Moralische Fragwürdigkeit
\item Wirtschaftliche Interessen
\item Fehlende Berücksichtigung der \\ gesellschaftlichen Folgen
\end{itemize}
}}}}

\makebox(0,450){}

\colorbox{gray!20}{
\makebox(900,290){
%right of Entwicklung
\large
\scalebox{1.2}{
\parbox{850pt}{
\begin{itemize}[$\bullet$]
\item Fehlende Dokumentation
\item Vernachlässigung des Datenschutzes
\end{itemize}
}}}}

\makebox(0,450){}

\colorbox{gray!20}{
\makebox(900,630){
%right of Anwendung
\large
\scalebox{1.2}{
\parbox{850pt}{
\begin{itemize}[$\bullet$]
\item Fehlende Nachvollziehbarkeit
\item Unvorhergesehene Einflüsse
\item Fehlendes Vertrauen in die Technik
\item Unvorhergesehene Folgen
\item Ungeklärte Haftung
\end{itemize}
}}}}

\makebox(900,350){
%bottom right
}

\end{minipage}


%TODO ask if this the following is ok if so change to iftrue or remove iffalse and fi
%creative commons with attribution and share-alike
\iffalse{

\hfill

\large
{
\color{white}
\llap{\makebox[16cm]{\ccLogo \ccAttribution \ccShareAlike}}
}
}
\fi

\clearpage

%un-comment to generate testpage
%\NoBgThispage
\layout
\clearpage


\end{document}